\documentclass[10pt,a4paper]{amsart}
\usepackage[margin=19mm]{geometry}
\usepackage{amsmath,amssymb,mathtools}
\usepackage[scr=rsfs,cal=euler]{mathalfa}
\usepackage{xfrac}
\usepackage[unicode,hidelinks,bookmarks=false]{hyperref}
\newcommand{\ie}{i.e.\ }
\newcommand{\cf}{cf.\ }
\newcommand{\resp}{respectively\ }
\newcommand{\Subsection}{\subsection}
\providecommand{\nobreakdash}{\nobreak\hbox{-}}
\setlength{\emergencystretch}{2em}
\multlinegap0pt
\usepackage{bm}
\usepackage{bbm}
\usepackage{dsfont}
\usepackage{xspace}
\usepackage{cancel}
\usepackage{mathtools}
\usepackage{amsthm}
\makeatletter
\@ifundefined{theorem}{\newtheorem{theorem}{Theorem}}{}
\@ifundefined{lemma}{\newtheorem{lemma}[theorem]{Lemma}}{}
\makeatother
\DeclarePairedDelimiter{\set}{\lbrace}{\rbrace}
\newcommand{\fitbracket}[1]{\left[#1\right]}
\newcommand{\fitparenth}[1]{\left(#1\right)}
\newcommand{\inprod}[1]{\left\langle#1\right\rangle}
\newcommand{\separation}{\Delta}
\newcommand{\vect}[1]{\mathbf{#1}}
\title[Diffusion models recover accurate mixture weights despite sc]{Diffusion models recover accurate mixture weights despite score function insensitivity}
\author{Andrew Dennehy, Ramchandran Muthukumar, Rebecca Willett, Nisha Chandramoorthy}
\date{Source: arXiv:2607.15485v1 (CC BY 4.0)}
\begin{document}
\maketitle

\noindent\emph{Source and licence.} Diffusion models recover accurate mixture weights despite score function insensitivity. Version arXiv:2607.15485v1 (CC BY 4.0), distributed under the Creative Commons Attribution 4.0 International licence, as recorded in the arXiv licence metadata for this exact version and shown on the abstract page https://arxiv.org/abs/2607.15485v1. Full rights notice: licenses/arxiv-2607.15485-CC-BY-4.0.txt.\par\smallskip

\noindent\textit{Editorial scope.} Counted unit: the Lemma whose source label is \texttt{lemma:bounds} in the article (source0.tex lines 1627--1650, source label \texttt{lemma:bounds}), delivered together with the article's own complete original proof of it (source0.tex lines 1652--1684, one \texttt{proof} environment of 33 lines). No mathematics is written, repaired or re-derived: statement and proof are the article's own text line for line. Required context retained uncounted: none. Every definition, display and label of those lines is retained unchanged. Omitted from this package and not redistributed: 1--1626, 1651--1651, 1685--2011. Nothing inside a retained excerpt is omitted or altered: every retained body is delivered exactly as the article's lines are written, so no omission receipt of any kind accompanies this package. Citations: the retained excerpts carry no citation command at all, so no bibliography entry of the article is transcribed and no reference is redistributed. Reference adaptation: none was needed and none was made; every reference inside the delivered unit resolves inside it, so no unresolved reference or citation is produced. Printed slips kept literal: no printed word, symbol, hypothesis or number of the counted statement or of its proof was altered. Layout and class: the article's own class is \texttt{article} with packages including amsmath, amssymb, amsfonts, amsthm, thmtools, natbib, xcolor, hyperref, xurl, cleveref, microtype, graphicx, subcaption, booktabs; the wrapper is the lane's pinned \texttt{amsart} editorial wrapper at 10pt on A4, which restates the theorem-like environments the retained text needs and adds the article's own macro definitions that the retained text needs (\texttt{\textbackslash fitbracket}, \texttt{\textbackslash fitparenth}, \texttt{\textbackslash inprod}, \texttt{\textbackslash separation}, \texttt{\textbackslash set}, \texttt{\textbackslash vect}), with extra wrapper packages (bm, bbm, dsfont, xspace, cancel, mathtools). The delivered numbering is the unit's own. Assets: no retained span contains a figure environment, a graphics-inclusion command, a PDF-page inclusion, a table or a TikZ picture; no image file is copied into this package, the package declares no asset dependency and no image file is redistributed with it. Licence: version arXiv:2607.15485v1 of this article is distributed under the Creative Commons Attribution 4.0 International licence, as recorded in the arXiv licence metadata for this exact version and shown on the abstract page https://arxiv.org/abs/2607.15485v1. A full rights notice is delivered at licenses/arxiv-2607.15485-CC-BY-4.0.txt. Characters: every retained excerpt is pure ASCII, so the delivered engine, XeLaTeX with its default Latin Modern text font, reports no missing character.
\begin{lemma}\label{lemma:bounds}
    Let $\vect z_t$ and $\vect z_0$ be as above. Let $M(\vect z_0):=\inprod{\vect z_0 - \mu_0, \mu_1 - \mu_0}$, and define the log-odds ratios 
    \begin{align*}
        \hat r &:= \log \frac{1-{\hat\gamma}}{{\hat\gamma}}
        \\
        r^* &:= \log \frac{1-\gamma^*}{\gamma^*}
    \end{align*}
    
    Let $\epsilon>0$ be such that 
    \begin{align*}
    \epsilon &\leq \exp\fitparenth{-\frac{1}{8}\fitbracket{(\hat r-r^*)^2 + \fitparenth{\hat r+r^*-\frac{M(\vect z_0)^2}{\separation^2}}_+^2}}
    \end{align*}
    where $(\cdot)_+:= \max\set{\cdot, 0}$.
    
    Let
    \begin{align*}
        a&=\max\left\{\frac{M(\vect z_0)}{\separation^2} + \frac{1}{\separation}\sqrt{\fitparenth{\frac{M(\vect z_0)^2}{\separation^2} - \fitparenth{ \hat r + r^*}-\sqrt{8\log \frac{1}{\epsilon}-\fitparenth{\hat r-r^*}^2}}_+},\sqrt{\alpha_1}\right\}
        \\
        b&= \min\left\{\frac{M(\vect z_0)}{\separation^2} + \frac{1}{\separation}\sqrt{\frac{M(\vect z_0)^2}{\separation^2} - \fitparenth{ \hat r + r^*}+\sqrt{8\log \frac{1}{\epsilon}-\fitparenth{\hat r-r^*}^2}},1\right\}
    \end{align*}

    Then for all $t$ such that $\sqrt{\alpha_t}\in [a,b]$, 
    \[e^{-\frac{1}{4}\fitparenth{f_0(\vect z_0,t)-\hat r }^s2-\frac{1}{4}\fitparenth{f_0(\vect z_0,t)-r^*}^2} \geq \epsilon\]
\end{lemma}

\begin{proof}
    First, the inequality is equivalent to 
    \begin{align*}
       \fitparenth{f_0(\vect z_0,t)-r^*}^2+\fitparenth{f_0(z_0,t)-r^*}^2\leq 4\log \frac{1}{\epsilon}
    \end{align*}

    
    
    
    Plugging in the definition of $f$, the inequality holds whenever 
    \[\fitbracket{\frac{\separation^2}{2} \fitparenth{\sqrt{\alpha_t}-\frac{M(\vect z_0)}{\separation^2}}^2 - \frac{M(\vect z_0)^2}{2\separation^2}+ \frac{\hat r+r^*}{2}}^2 \leq 2\log \frac{1}{\epsilon}-\frac{(\hat r-r^*)^2}{4} \]
    
    We wish to identify the values of $\sqrt{\alpha_t}$ which make this inequality hold. This function is continuous, so its sign can only change at the roots of the corresponding equality constraint:
    \[\fitbracket{\frac{\separation^2}{2} \fitparenth{\sqrt{\alpha_t}-\frac{M(\vect z_0)}{\separation^2}}^2 - \frac{M(\vect z_0)^2}{2\separation^2}+ \frac{\hat r+r^*}{2}}^2 = 2\log \frac{1}{\epsilon}-\frac{(\hat r -r^*)^2}{4}\]

    Solving this quartic gives the roots
    \[\sqrt{\alpha_t}=\frac{M(\vect z_0)}{\separation^2} \pm\sqrt{\frac{M(\vect z_0)^2}{\separation^4}-\frac{\hat r+r^*}{\separation^2}\pm \frac{1}{\separation^2}\sqrt{8\log \frac{1}{\epsilon}-(\hat r-r^*)^2}}\]
    
    The condition on $\epsilon$ is such that the square roots are all well-defined. Let $b'$ be the largest of these roots (both $\pm$'s are set to $+$). As $\sqrt{\alpha_t}\rightarrow \infty$, we would have $f_0(\vect z_0,t)\rightarrow -\infty$ (because $-\separation^2/2$ is strictly negative), and hence 
    \[e^{-\frac{1}{4}\fitparenth{f_0(\vect z_0,t)-\hat r }^2-\frac{1}{4}\fitparenth{f_0(\vect z_0,t)-r^*}^2} \rightarrow 0\]
    
    Therefore, for $\sqrt{\alpha_t}>b'$ the inequality does not hold. Note that the four roots above are distinct (with two being complex for $\epsilon$ sufficiently small), meaning that each root has multiplicity 1, and hence the polynomial changes signs at $b'$. Therefore, for $\sqrt{\alpha_t}$ between $b'$ and the next largest real root, the inequality holds. The next largest real root is
    \[a'=\begin{cases}
        \frac{M(\vect z_0)}{\separation^2} +\sqrt{\frac{M(\vect z_0)^2}{\separation^4}-\frac{\hat r+r^*}{\separation^2}-\frac{1}{\separation^2}\sqrt{8\log \frac{1}{\epsilon}-(\hat r-r^*)^2}} & \text{if}\quad \frac{M(\vect z_0)^2}{\separation^2}-(\hat r+r^*)-\sqrt{8\log \frac{1}{\epsilon}-(\hat r-r^*)^2} \geq0
        \\
        \frac{M(\vect z_0)}{\separation^2} -\sqrt{\frac{M(\vect z_0)^2}{\separation^4}-\frac{\hat r+r^*}{\separation^2}+\frac{1}{\separation^2}\sqrt{8\log \frac{1}{\epsilon}-(\hat r-r^*)^2}} & \text{otherwise}
    \end{cases}\]
    
    In the second case we have $a'\leq \frac{M(\vect z_0)}{\separation^2}$, so
    \[a'\leq \frac{M(\vect z_0)}{\separation^2} +\sqrt{\fitparenth{\frac{M(\vect z_0)^2}{\separation^4}-\frac{\hat r+r^*}{\separation^2}-\frac{1}{\separation^2}\sqrt{8\log \frac{1}{\epsilon}-(\hat r-r^*)^2}}_+}=:a'' \]
    
    Therefore, the inequality holds on $[a',b']$ and $[a'',b']\subseteq [a',b']$, so the inequality holds on $[a'',b']$. Note that for $t\in [0,1]$, $\alpha_t\in [\alpha_1, 1]$ where $\alpha_1 >0$ is small, so we let $b:= \min\set{b', 1}$ and $a:=\max\set{a,\sqrt{\alpha_1}}$.
\end{proof}

\end{document}
